\chapter{Computational Geometry}

Up to now, the objects our algorithms have chewed on were numbers, characters,
keys in a tree. Comparisons were cheap: is $5 < 8$? is the letter \texttt{a}
before \texttt{b}? But the world is not made of numbers in a line. It is made of
\emph{places} -- pixels on a screen, cities on a map, atoms in a protein, robots
on a factory floor. Each of those is a point in space, and the questions we want
to ask about them are geometric: which two points are closest? does this road
cross that river? what is the smallest fence that encloses this herd?

These questions \emph{feel} obvious to a human eye. You can glance at a scatter
of dots and instantly see the two that are nearest, or trace the outline that
wraps them all. Your eye is doing something an algorithm cannot do for free: it
is processing the whole picture at once. Our job in this chapter is to teach a
machine -- which sees only a list of coordinate pairs, one at a time -- to answer
the same questions, and to do it \emph{fast}. The recurring miracle will be that
a tiny algebraic gadget, the \emph{cross product}, replaces all that visual
intuition with a single arithmetic sign.

We will build four classic results from the ground up:
\begin{itemize}
  \item the \emph{orientation test} -- given three points, do they turn left,
        turn right, or lie on a line? This is the atom from which everything else
        is built.
  \item \emph{segment intersection} -- do two line segments cross, using only
        orientation tests and no division at all.
  \item the \emph{convex hull} -- the smallest convex polygon enclosing a set of
        points -- via two beautiful algorithms, \emph{Graham's scan} and
        \emph{Jarvis's march}.
  \item the \emph{closest pair of points} -- found in $\Theta(n \log n)$ by a
        divide-and-conquer trick that looks, at first, like it cannot possibly
        work.
\end{itemize}

\begin{historybox}
This chapter maps to CLRS Chapter~33, ``Computational Geometry.''
Section~33.1 is line-segment properties (our orientation test and segment
intersection), 33.2 is the sweep for ``do any segments intersect,'' 33.3 is
convex hull (Graham's scan and Jarvis's march), and 33.4 is the closest pair.
The field itself was crystallized in the 1970s and 1980s -- Graham's scan dates
to a 1972 paper by Ronald Graham motivated, charmingly, by a problem at Bell
Labs; the closest-pair divide-and-conquer is due to Shamos and Hoey. If you ever
take a dedicated course, the canonical text is de~Berg, Cheong, van~Kreveld, and
Overmars, \emph{Computational Geometry: Algorithms and Applications}.
\end{historybox}

\chapterline

\section{The cross product: the only tool you really need}

Everything starts with one question about \emph{three} points. Let us name them
$p_0$, $p_1$, $p_2$. Imagine you walk from $p_0$ straight to $p_1$, then you must
turn to head toward $p_2$. Do you turn \emph{left}, do you turn \emph{right}, or
do you keep going \emph{straight}? That single yes/no/no answer -- which way you
turned -- is the seed of this entire chapter.

\begin{intuitionbox}
You are driving. You reach $p_1$ and you want to get to $p_2$. If the steering
wheel goes counterclockwise (left), the three points make a left turn. If it goes
clockwise (right), a right turn. If you do not touch the wheel at all, the three
points are \emph{collinear} -- they sit on one straight line. ``Left, right, or
straight'' is the whole vocabulary. Astonishingly, almost every geometric
algorithm in this chapter is built by asking this question over and over.
\end{intuitionbox}

How do we compute the turn direction from raw coordinates, without a steering
wheel? We use \emph{vectors} and the \emph{cross product}. Treat each point as a
pair $p = (p_x, p_y)$. From $p_0$ form two arrows:
\[
  \vec{u} = p_1 - p_0 = (p_{1x} - p_{0x},\; p_{1y} - p_{0y}),
  \qquad
  \vec{v} = p_2 - p_0 = (p_{2x} - p_{0x},\; p_{2y} - p_{0y}).
\]
The \emph{two-dimensional cross product} of $\vec{u}$ and $\vec{v}$ is the single
number
\[
  \vec{u} \times \vec{v} \;=\; u_x v_y - u_y v_x .
\]
That is it -- two multiplications and a subtraction. The magic is in its
\emph{sign}.

\begin{formalbox}
Let $\vec{u} = (u_x, u_y)$ and $\vec{v} = (v_x, v_y)$. Define
$\vec{u} \times \vec{v} = u_x v_y - u_y v_x$. Then:
\begin{itemize}
  \item if $\vec{u} \times \vec{v} > 0$, then $\vec{v}$ is counterclockwise from
        $\vec{u}$ (a left turn);
  \item if $\vec{u} \times \vec{v} < 0$, then $\vec{v}$ is clockwise from
        $\vec{u}$ (a right turn);
  \item if $\vec{u} \times \vec{v} = 0$, then $\vec{u}$ and $\vec{v}$ are
        \emph{parallel} (collinear through $p_0$).
\end{itemize}
The magnitude $\abs{\vec{u} \times \vec{v}}$ equals the area of the parallelogram
spanned by $\vec{u}$ and $\vec{v}$, so half of it is the area of the triangle
$p_0 p_1 p_2$.
\end{formalbox}

Why does the sign encode the turn? Here is the gentle version. The cross product
$u_x v_y - u_y v_x$ is a \emph{signed area}. If you sweep the arrow $\vec{u}$
counterclockwise toward $\vec{v}$, the swept region has positive area; sweep it
clockwise and you get negative area; if there is nothing to sweep because they
point the same way, the area is zero. The sign of an area tells you the
\emph{direction of rotation}. That is the entire trick, and we never needed
trigonometry, angles, or square roots -- just multiply and subtract.

\begin{focusbox}
The cross product replaces ``looking at the picture.'' Where your eye sees a left
turn, the machine computes a positive number. Hold this equivalence in your head:
\[
  \text{sign of } (u_x v_y - u_y v_x)
  \;\;\longleftrightarrow\;\;
  \text{left / straight / right.}
\]
Everything else in this chapter is built by asking this question, sometimes
thousands of times.
\end{focusbox}

\subsection{The \textsc{Orientation} primitive}

Let us bottle this up as a reusable function. Given three points it returns one
of three labels. We will encode them as $+1$ (counterclockwise / left), $-1$
(clockwise / right), and $0$ (collinear).

\begin{algorithm}[H]
\caption{\textsc{Orientation}($p_0, p_1, p_2$)}
\KwIn{Three points $p_0, p_1, p_2$ in the plane.}
\KwOut{$+1$ if $p_0 p_1 p_2$ make a counterclockwise (left) turn, $-1$ if clockwise (right), $0$ if collinear.}
$c \leftarrow (p_{1x} - p_{0x})(p_{2y} - p_{0y}) - (p_{1y} - p_{0y})(p_{2x} - p_{0x})$\;
\If{$c > 0$}{\Return $+1$}
\If{$c < 0$}{\Return $-1$}
\Return $0$\;
\end{algorithm}

Let us trace it on real numbers. Take $p_0 = (0,0)$, $p_1 = (4,0)$, and three
candidates for $p_2$.

\begin{example}\label{ch32:orient-trace}
With $p_0 = (0,0)$ and $p_1 = (4,0)$ (so $\vec{u} = (4,0)$, pointing east):
\begin{itemize}
  \item $p_2 = (2, 3)$ (north of the line): $c = (4)(3) - (0)(2) = 12 > 0$, a
        \textbf{left turn}. Yes -- to get from heading east to pointing at a spot
        above the line, you turn left. \checkmark
  \item $p_2 = (2, -3)$ (south of the line): $c = (4)(-3) - (0)(2) = -12 < 0$, a
        \textbf{right turn}. \checkmark
  \item $p_2 = (7, 0)$ (straight ahead): $c = (4)(0) - (0)(7) = 0$,
        \textbf{collinear}. \checkmark
\end{itemize}
\end{example}

Here is the picture the algebra of Example~\ref{ch32:orient-trace} is describing.

\begin{center}
\begin{tikzpicture}[scale=0.9,>=Stealth]
  \coordinate (p0) at (0,0);
  \coordinate (p1) at (4,0);
  \coordinate (pL) at (2,2);
  \coordinate (pR) at (2,-2);
  \draw[->,thick,sourceblue] (p0) -- (p1) node[midway,below=2pt]{$\vec u$};
  \draw[->,thick,targetgreen] (p0) -- (pL) node[midway,above left=-2pt]{left ($c>0$)};
  \draw[->,thick,trapred] (p0) -- (pR) node[midway,below left=-2pt]{right ($c<0$)};
  \draw[dashed,gray] (p1) -- (6,0) node[right]{collinear ($c=0$)};
  \fill (p0) circle (2pt) node[above left]{$p_0$};
  \fill (p1) circle (2pt) node[below right]{$p_1$};
  \fill (pL) circle (2pt) node[above]{$p_2$ (left)};
  \fill (pR) circle (2pt) node[below]{$p_2$ (right)};
\end{tikzpicture}
\end{center}

\begin{trapbox}
Two traps bite beginners here. First, \textbf{do not compute the slope}
$(p_{1y}-p_{0y})/(p_{1x}-p_{0x})$ to compare directions. Division invites
division-by-zero (vertical lines!) and floating-point noise. The cross product
uses only $\times$ and $-$, so on \emph{integer} inputs it is \emph{exact} -- no
rounding, ever. Second, \textbf{order matters}: \textsc{Orientation}$(p_0,p_1,p_2)$
and \textsc{Orientation}$(p_0,p_2,p_1)$ have opposite signs. Swapping the last two
arguments flips left and right, because you are now asking about the reverse turn.
\end{trapbox}

\begin{aestheticbox}
This one sign computation is everywhere. A GPS deciding which side of a road you
are on, a CAD program checking whether you clicked inside a polygon, a video game
testing whether a bullet's path crosses a wall, a chip-design tool checking that
wires do not overlap -- all of them lean on the orientation test. It is the
``\texttt{<}'' of the geometric world. Master it once and you have the keystone.
\end{aestheticbox}

\section{Do two segments intersect?}

Now a real question. You have two line segments, $\overline{p_1 p_2}$ and
$\overline{p_3 p_4}$. Do they cross? Not the infinite lines through them -- the
finite segments, with their endpoints. This is the test a game engine runs to see
if your character walked through a wall.

The naive approach is to solve for the intersection point of the two lines with
algebra, then check whether that point falls inside both segments. It works, but
it needs division, it stumbles on parallel and vertical lines, and floating-point
error makes the ``is the point inside?'' check shaky. We will do something
cleaner that uses \emph{only orientation tests} -- exact integer arithmetic, no
division.

\begin{intuitionbox}
Here is the key idea, and it is lovely. Segment $\overline{p_1 p_2}$ crosses
segment $\overline{p_3 p_4}$ exactly when each segment ``straddles'' the line
through the other. ``Straddle'' means the two endpoints sit on \emph{opposite
sides}. Look at $p_3$ and $p_4$: if one is to the left of the line $p_1 p_2$ and
the other is to the right, then $\overline{p_3 p_4}$ pokes across that line. If
\emph{both} segments straddle each other's lines, they must cross. If either one
fails to straddle, they miss.
\end{intuitionbox}

``Which side of the line $p_1 p_2$ is the point $p_3$ on?'' is precisely
\textsc{Orientation}$(p_1, p_2, p_3)$. So the whole test is four orientation
calls. Let us define
\[
  d_1 = \textsc{Orientation}(p_3, p_4, p_1), \quad
  d_2 = \textsc{Orientation}(p_3, p_4, p_2),
\]
\[
  d_3 = \textsc{Orientation}(p_1, p_2, p_3), \quad
  d_4 = \textsc{Orientation}(p_1, p_2, p_4).
\]
The points $p_1$ and $p_2$ are on opposite sides of line $p_3 p_4$ exactly when
$d_1$ and $d_2$ have \emph{opposite signs}; symmetrically for $d_3, d_4$ and the
line $p_1 p_2$.

\begin{center}
\begin{tikzpicture}[scale=0.85,>=Stealth]
  % crossing case
  \begin{scope}
    \coordinate (a) at (0,0); \coordinate (b) at (4,1);
    \coordinate (c) at (1,2); \coordinate (d) at (3,-1);
    \draw[thick,sourceblue] (a)--(b);
    \draw[thick,trapred] (c)--(d);
    \fill (a) circle(1.6pt) node[left]{$p_1$};
    \fill (b) circle(1.6pt) node[right]{$p_2$};
    \fill (c) circle(1.6pt) node[above]{$p_3$};
    \fill (d) circle(1.6pt) node[below]{$p_4$};
    \node at (2,-2){they cross};
  \end{scope}
  % non-crossing case
  \begin{scope}[xshift=7cm]
    \coordinate (a) at (0,0); \coordinate (b) at (4,1);
    \coordinate (c) at (1,2); \coordinate (d) at (3,2.5);
    \draw[thick,sourceblue] (a)--(b);
    \draw[thick,trapred] (c)--(d);
    \fill (a) circle(1.6pt) node[left]{$p_1$};
    \fill (b) circle(1.6pt) node[right]{$p_2$};
    \fill (c) circle(1.6pt) node[above]{$p_3$};
    \fill (d) circle(1.6pt) node[above]{$p_4$};
    \node at (2,-2){no crossing: $p_3,p_4$ same side};
  \end{scope}
\end{tikzpicture}
\end{center}

There is one subtlety: what if an endpoint lands \emph{exactly} on the other
segment? Then some $d_i = 0$, and we must check whether that collinear point
actually lies \emph{within} the bounds of the segment (not out on the extension).
A tiny helper, \textsc{OnSegment}, handles it: assuming $p_i$ is known collinear
with $\overline{p_j p_k}$, it just checks that $p_i$'s coordinates fall inside the
bounding box of that segment.

\begin{algorithm}[H]
\caption{\textsc{SegmentsIntersect}($p_1, p_2, p_3, p_4$)}
\KwIn{Four points defining segments $\overline{p_1 p_2}$ and $\overline{p_3 p_4}$.}
\KwOut{\textsc{true} if the segments intersect (including touching), else \textsc{false}.}
$d_1 \leftarrow \textsc{Orientation}(p_3,p_4,p_1)$\;
$d_2 \leftarrow \textsc{Orientation}(p_3,p_4,p_2)$\;
$d_3 \leftarrow \textsc{Orientation}(p_1,p_2,p_3)$\;
$d_4 \leftarrow \textsc{Orientation}(p_1,p_2,p_4)$\;
\If{$d_1 d_2 < 0$ \textbf{and} $d_3 d_4 < 0$}{\Return \textsc{true} \tcp*{proper crossing}}
\If{$d_1 = 0$ \textbf{and} \textsc{OnSegment}($p_3,p_4,p_1$)}{\Return \textsc{true}}
\If{$d_2 = 0$ \textbf{and} \textsc{OnSegment}($p_3,p_4,p_2$)}{\Return \textsc{true}}
\If{$d_3 = 0$ \textbf{and} \textsc{OnSegment}($p_1,p_2,p_3$)}{\Return \textsc{true}}
\If{$d_4 = 0$ \textbf{and} \textsc{OnSegment}($p_1,p_2,p_4$)}{\Return \textsc{true}}
\Return \textsc{false}\;
\end{algorithm}

Example~\ref{ch32:seg-trace} works the four orientation calls end to end. Here the
products $d_1 d_2 < 0$ and $d_3 d_4 < 0$ are a slick way to say ``the two
signs differ.'' If they differ \emph{and} neither is zero, the segments straddle
properly and cross. The remaining four lines handle the boundary cases where an
endpoint touches.

\begin{example}\label{ch32:seg-trace}
Take $\overline{p_1 p_2}$ with $p_1=(0,0)$, $p_2=(4,1)$, and $\overline{p_3 p_4}$
with $p_3=(1,2)$, $p_4=(3,-1)$ (the crossing picture above). Compute:
\begin{itemize}
  \item $d_1=\textsc{Orientation}(p_3,p_4,p_1)$: $\vec u = p_4-p_3=(2,-3)$,
        $\vec v = p_1-p_3=(-1,-2)$, $c=(2)(-2)-(-3)(-1)=-4-3=-7<0$, so $d_1=-1$.
  \item $d_2=\textsc{Orientation}(p_3,p_4,p_2)$: $\vec v=p_2-p_3=(3,-1)$,
        $c=(2)(-1)-(-3)(3)=-2+9=7>0$, so $d_2=+1$.
  \item $d_1 d_2 = -1 < 0$: $p_1,p_2$ straddle line $p_3 p_4$. \checkmark
\end{itemize}
$d_3=\textsc{Orientation}(p_1,p_2,p_3)$: $\vec u=(4,1)$, $\vec v=(1,2)$,
$c=8-1=7>0$, so $d_3=+1$. And $d_4$: $\vec v=p_4-p_1=(3,-1)$, $c=(4)(-1)-(1)(3)
=-7<0$, so $d_4=-1$. Then $d_3 d_4=-1<0$. Both straddles hold, so the segments
\textbf{intersect}. \checkmark
\end{example}

\begin{complexitybox}
\textsc{SegmentsIntersect} runs in $\Th(1)$ time: a fixed number of orientation
tests, each $\Th(1)$. On integer coordinates it is \emph{exact} -- it never
rounds. Now suppose instead you have $n$ segments and ask ``do \emph{any} two of
them intersect?'' Checking all pairs is $\Th(n^2)$. CLRS~33.2 shows a
\emph{sweep-line} algorithm that answers this yes/no question in $\Th(n \log n)$
by sweeping a vertical line left to right and keeping the segments it currently
crosses in a balanced BST ordered by height -- two segments can only cross if
they are ever \emph{adjacent} in that order, so you only test neighbors.
\end{complexitybox}

\section{The convex hull}

Now the headline problem. You are given $n$ points scattered in the plane. The
\emph{convex hull} is the smallest convex polygon that contains all of them.

\begin{intuitionbox}
Hammer a nail into the table at each point. Stretch a rubber band wide open around
the whole cluster and let go. It snaps tight and clings to the outermost nails,
ignoring everything tucked inside. The taut rubber band is the convex hull. ``No
dents'': between any two band-nails the band goes perfectly straight, never caving
inward. Equivalently, the hull is convex -- the segment between any two points of
the polygon stays inside it.
\end{intuitionbox}

\begin{center}
\begin{tikzpicture}[scale=0.8]
  \foreach \p in {(0,0),(4,0.4),(5,3),(2.5,4),(-0.5,2.8),(2,2),(3,1.5),(1.5,1),(3.5,2.6)}
    \fill[gray] \p circle (2.2pt);
  \draw[thick,sourceblue] (0,0)--(4,0.4)--(5,3)--(2.5,4)--(-0.5,2.8)--cycle;
  \node[sourceblue] at (5.6,1.5){hull};
  \node[gray] at (2.5,1.6){interior points};
\end{tikzpicture}
\end{center}

The hull points are exactly the ``extreme'' ones; every other point is interior
and irrelevant to the boundary. We will see two algorithms. Both lean entirely on
the orientation test.

\begin{aestheticbox}
Convex hulls are everywhere a ``smallest enclosing shape'' is useful: collision
detection in physics engines and games (it is cheap to test convex shapes),
finding the diameter of a point set, path planning for robots (hug the hull to
avoid obstacles), pattern recognition, and as a preprocessing step for dozens of
other geometric algorithms. When you ask a phone to crop to ``the region with all
the faces,'' a hull is lurking nearby.
\end{aestheticbox}

\subsection{Graham's scan}

Graham's scan sorts the points by angle around an anchor, then walks through them
keeping a ``stack'' of hull vertices, popping any point that would create a
right turn (a dent). It runs in $\Th(n \log n)$, dominated by the sort.

\textbf{Step 1 -- pick the anchor.} Choose $p_0$, the point with the lowest
$y$-coordinate (break ties by lowest $x$). This point is \emph{guaranteed} to be
on the hull -- nothing sits below it, so the rubber band must touch it.

\textbf{Step 2 -- sort by polar angle.} Sort the remaining points by the angle
the ray $p_0 \to p$ makes with the positive $x$-axis, increasing. Here is the
elegant part: you do \emph{not} compute any angles. To decide whether $p_a$ comes
before $p_b$, just ask \textsc{Orientation}$(p_0, p_a, p_b)$. If it is $+1$
(counterclockwise), $p_a$ has the smaller angle. The cross product sorts by angle
\emph{without ever calling a trig function}. (Ties at equal angle: keep the
farther point, or handle collinear points with a small rule.)

\textbf{Step 3 -- the scan.} Now walk the sorted points, maintaining a stack of
the hull-so-far. Push $p_0$ and the first two sorted points. Then for each
subsequent point $p_i$: while the top two stack entries together with $p_i$ make a
\emph{non-left} turn (right turn or collinear), \emph{pop} the top -- it was a
dent, a mistake we now correct. Then push $p_i$. When you have consumed every
point, the stack \emph{is} the hull, in counterclockwise order.

\begin{algorithm}[H]
\caption{\textsc{GrahamScan}($P$)}
\KwIn{A set $P$ of $n \ge 3$ points.}
\KwOut{The vertices of the convex hull of $P$ in counterclockwise order.}
$p_0 \leftarrow$ point of $P$ with lowest $y$ (ties: lowest $x$)\;
sort the other points by polar angle around $p_0$, using \textsc{Orientation} as the comparator; call them $p_1, \dots, p_{m}$\;
$S \leftarrow$ empty stack\;
\textsc{Push}($S, p_0$);\; \textsc{Push}($S, p_1$);\; \textsc{Push}($S, p_2$)\;
\For{$i \leftarrow 3$ \KwTo $m$}{
  \While{$\abs{S} \ge 2$ \textbf{and} \textsc{Orientation}(\textsc{NextToTop}($S$), \textsc{Top}($S$), $p_i$) $\le 0$}{
    \textsc{Pop}($S$) \tcp*{the turn is not left: pop the dent}
  }
  \textsc{Push}($S, p_i$)\;
}
\Return $S$\;
\end{algorithm}

\begin{labbox}
Let us run it on a small set. Anchor $p_0=(0,0)$, and after sorting by angle
suppose the points are
\[
  p_1=(3,1),\; p_2=(5,3),\; p_3=(2,2),\; p_4=(0,4).
\]
Start the stack as $[\,p_0, p_1, p_2\,]$ (bottom to top).
\begin{itemize}
  \item \textbf{Consider $p_3=(2,2)$.} Check the turn $p_1 \to p_2 \to p_3$:
        $\vec u=p_2-p_1=(2,2)$, $\vec v=p_3-p_1=(-1,1)$, $c=(2)(1)-(2)(-1)=4>0$,
        a \emph{left} turn. Left is good -- no pop. Push $p_3$. Stack:
        $[p_0,p_1,p_2,p_3]$. (At this moment $p_3$ looks like a hull vertex; the
        next point will reveal that it is not.)
  \item \textbf{Consider $p_4=(0,4)$.} Check $p_2 \to p_3 \to p_4$:
        $\vec u=p_3-p_2=(-3,-1)$, $\vec v=p_4-p_2=(-5,1)$,
        $c=(-3)(1)-(-1)(-5)=-3-5=-8<0$, a \emph{right} turn. Right is a dent --
        \emph{pop} $p_3$. Stack: $[p_0,p_1,p_2]$.
  \item Re-check $p_1 \to p_2 \to p_4$: $\vec u=p_2-p_1=(2,2)$,
        $\vec v=p_4-p_1=(-3,3)$, $c=(2)(3)-(2)(-3)=6+6=12>0$, a left turn. Good --
        push $p_4$. Stack: $[p_0,p_1,p_2,p_4]$.
\end{itemize}
The hull is $p_0,p_1,p_2,p_4 = (0,0),(3,1),(5,3),(0,4)$, and $p_3=(2,2)$ was
correctly discarded as interior. Trace this on a walk and watch the dent get
popped.
\end{labbox}

Why does popping on a non-left turn give the right answer? Here is the invariant.

\begin{focusbox}
\textbf{Loop invariant (Graham's scan).} At the top of each \texttt{for}
iteration, the stack $S$ holds, from bottom to top, the vertices of the convex
hull of the points considered so far, in counterclockwise order, with every
consecutive triple making a left turn.
\end{focusbox}

\emph{Initialization.} Three points are trivially their own (degenerate) hull.
\emph{Maintenance.} When we add $p_i$, the \texttt{while} loop pops every vertex
that would create a right or straight turn -- exactly the vertices that $p_i$ now
``hides'' from the boundary, the dents. After popping, the top triple turns left,
so the invariant is restored when we push $p_i$. The points sit in angle order, so
$p_i$ is the next counterclockwise vertex. \emph{Termination.} After the last
point, $S$ is the hull of all of $P$. A clean induction, and the engine of the
proof is, again, the orientation sign.

\begin{complexitybox}
\textbf{Graham's scan.} The sort costs $\Th(n \log n)$. The scan looks like it
could be $\Th(n^2)$ because of the inner \texttt{while}, but each point is
\emph{pushed once and popped at most once} over the whole run -- an amortized
argument (Chapter~17). So the scan itself is $\Th(n)$, and the total is
\[
  T(n) = \underbrace{\Th(n \log n)}_{\text{sort}} + \underbrace{\Th(n)}_{\text{scan}}
       = \Th(n \log n).
\]
The sort dominates. This is asymptotically optimal for comparison-based hull
algorithms.
\end{complexitybox}

\begin{trapbox}
Two classic mistakes. (1) Forgetting to pop in a \texttt{while} loop -- one pop is
not enough; adding a far-out point can invalidate \emph{several} previous
vertices, so you must keep popping until the turn is left. (2) Mishandling
collinear points. With $\le 0$ in the test, collinear points ($c=0$) get popped,
which usually you want (a vertex in the middle of a straight edge is not a true
corner). If your application needs them \emph{kept}, change the test to $<0$. Be
deliberate about which you mean.
\end{trapbox}

\subsection{Jarvis's march (gift wrapping)}

There is a second algorithm that is slower in the worst case but irresistibly
intuitive, and faster when the hull has few vertices. It is called \emph{Jarvis's
march}, or \emph{gift wrapping}, and it does exactly what your hand does when you
wrap a present: it walks around the outside, always turning as little as possible.

\begin{intuitionbox}
Start at the lowest point -- definitely on the hull. Now imagine sweeping a ray
counterclockwise until it bumps into the very first point. That point is the next
hull vertex; the ray ``wrapped'' onto it. Stand there and sweep again, and again,
each time pivoting to the next point, until you come all the way back to the
start. You have wrapped a string around the outside of the point cloud.
\end{intuitionbox}

\begin{center}
\begin{tikzpicture}[scale=0.8,>=Stealth]
  \coordinate (v) at (0,0);
  \coordinate (q) at (4,0.4);
  \foreach \p in {(5,3),(2.5,4),(-0.5,2.8),(2,2),(3,1.5),(1.5,1),(3.5,2.6)}
    \fill[gray] \p circle (2pt);
  \fill (v) circle (2.4pt) node[below left]{$v$ (current)};
  \fill (q) circle (2.4pt) node[below right]{$q$ (next vertex)};
  % the wrapping ray, and a faint cone showing everything is to its right
  \draw[->,thick,targetgreen] (v) -- (q);
  \draw[dashed,sourceblue] (v) -- (5,3);
  \draw[dashed,sourceblue] (v) -- (-0.5,2.8);
  \node[targetgreen] at (2,-0.55){every other point lies to the right of $v\!\to\!q$};
\end{tikzpicture}
\end{center}

The dashed rays in the figure show $v$ ``seeing'' the rest of the cloud; the solid
green ray $v \to q$ is the most clockwise of all of them, so $q$ is the next hull
vertex.

How do you find ``the next point'' without sweeping a physical ray? With -- you
guessed it -- the orientation test. From the current hull vertex $v$, you want the
point $q$ such that \emph{every other point} lies to the right of the ray
$v \to q$. Equivalently, $q$ is the most clockwise point as seen from $v$. You find
it by a linear scan: start with any candidate and replace it whenever
\textsc{Orientation}$(v, \text{candidate}, r) = -1$ (point $r$ is even more
clockwise).

\begin{algorithm}[H]
\caption{\textsc{JarvisMarch}($P$)}
\KwIn{A set $P$ of $n \ge 3$ points.}
\KwOut{The convex hull of $P$ in counterclockwise order.}
$\mathit{start} \leftarrow$ leftmost point of $P$ (ties: lowest $y$) \tcp*{guaranteed on the hull}
$v \leftarrow \mathit{start}$\;
$H \leftarrow$ empty list\;
\Repeat{$v = \mathit{start}$}{
  append $v$ to $H$\;
  $q \leftarrow$ any point of $P$ with $q \ne v$\;
  \For{each point $r \in P$}{
    \If{$\textsc{Orientation}(v, q, r) = -1$}{
      $q \leftarrow r$ \tcp*{$r$ is more clockwise; better next vertex}
    }
  }
  $v \leftarrow q$\;
}
\Return $H$\;
\end{algorithm}

\begin{complexitybox}
\textbf{Jarvis's march.} Each pass around the inner \texttt{for} loop is $\Th(n)$
and produces exactly one hull vertex. If the hull has $h$ vertices, the total cost
is $\Th(nh)$ -- this is an \emph{output-sensitive} bound. If only a handful of
points are on the hull ($h$ small, say a constant), Jarvis runs in $\Th(n)$,
beating Graham. But in the worst case \emph{all} $n$ points lie on the hull (think
of points on a circle), so $h = n$ and Jarvis degrades to $\Th(n^2)$. Graham's
$\Th(n \log n)$ is the safer default; Jarvis shines when you know the hull is
small.
\end{complexitybox}

\begin{walkbox}
Take a coin out of your pocket and put points only on its rim -- a perfect circle.
How many points are on the convex hull? \emph{All of them.} Now you understand
why Jarvis's march hits its $\Th(n^2)$ worst case here: $h = n$. Next, scatter $n$
points uniformly inside a square. Roughly how many end up on the hull? (The answer
grows like $\log n$, which is why Jarvis is often surprisingly fast in practice.)
\end{walkbox}

\section{The closest pair of points}

Final problem, and the prettiest. Given $n$ points, find the two whose Euclidean
distance is smallest. A radar system might run this to spot the two aircraft in
greatest danger of colliding.

The brute force is obvious: compute all $\binom{n}{2}$ pairwise distances and keep
the minimum. That is $\Th(n^2)$. Can we do better? Sorting was $\Th(n \log n)$;
surely distances are harder... and yet, with divide and conquer (Chapter~4), we
will hit $\Th(n \log n)$. The first time you see how, it feels like a magic trick.

\begin{intuitionbox}
Split the points into a left half and a right half with a vertical line. Find the
closest pair on the left (recursively) and the closest pair on the right
(recursively). Let $\delta$ be the smaller of those two distances. The true
closest pair is either entirely on the left, entirely on the right, \emph{or} it
straddles the line -- one point on each side. The first two cases are already
covered by $\delta$. The whole difficulty, and the whole cleverness, is the third:
checking straddling pairs \emph{fast}.
\end{intuitionbox}

\subsection{The setup and the strip}

First, presort: make a copy of the points sorted by $x$ (call it $X$) and another
sorted by $y$ (call it $Y$). We do this \emph{once}, up front, in $\Th(n \log n)$,
and then keep the sorted orders consistent through the recursion. Now the
recursion. Draw a vertical line at the median $x$, splitting the points into left
half $L$ and right half $R$ of size $n/2$ each. Recurse to get $\delta_L$ and
$\delta_R$, and set $\delta = \min(\delta_L, \delta_R)$.

Now for straddling pairs. A straddling pair closer than $\delta$ must have both
points within horizontal distance $\delta$ of the dividing line -- otherwise they
are already more than $\delta$ apart in $x$ alone. So we only care about points in
a vertical \emph{strip} of width $2\delta$ centered on the line.

\begin{center}
\begin{tikzpicture}[scale=0.78]
  \draw[->] (-0.3,0)--(8.3,0) node[right]{$x$};
  \draw[->] (0,-0.3)--(0,5) node[above]{$y$};
  \draw[dashed,thick,structurepurple] (4,-0.2)--(4,5) node[above]{median line};
  \fill[structurepurple!12] (3,-0.2) rectangle (5,5);
  \node[structurepurple] at (4,-0.6){width $2\delta$};
  % left points
  \foreach \p in {(0.6,1),(1.4,3.2),(2.2,0.6),(3.2,2.1),(3.6,3.8),(1.0,4.2)}
    \fill[sourceblue] \p circle (2pt);
  % right points
  \foreach \p in {(4.4,1.3),(4.6,3.4),(5.4,0.9),(6.4,2.6),(7.2,4),(5.0,4.4)}
    \fill[targetgreen] \p circle (2pt);
  \node[sourceblue] at (1.4,4.7){$L$};
  \node[targetgreen] at (6.6,4.7){$R$};
\end{tikzpicture}
\end{center}

Here is the worry. The strip could contain \emph{all} $n$ points (if they all
cluster near the line). Comparing every pair inside the strip would be $\Th(n^2)$
again, and the whole scheme collapses. We need the strip to be cheap.

\subsection{The seven-neighbors miracle}

This is the heart of the algorithm, and it is genuinely beautiful.

\begin{focusbox}
\textbf{The key geometric fact.} Sort the strip points by their $y$-coordinate.
Then for any point in the strip, you only need to compare it against the next
\emph{7} points in that $y$-order. Any straddling pair closer than $\delta$ must
appear within 7 positions of each other in the $y$-sorted strip.
\end{focusbox}

Why \emph{7}? Picture the strip as a $2\delta$-wide vertical band. Take any point
$p$ in it. Any point that could be within $\delta$ of $p$ must lie in a
$2\delta \times \delta$ rectangle just above $p$ (within $\delta$ in $y$, and at
most $2\delta$ across in $x$, with the median line running down its middle). Now
split that rectangle into a grid of squares of side $\delta/2$: it fits exactly
$4 \times 2 = 8$ such squares, four entirely on the left side of the line and four
entirely on the right. And here is the punchline: \emph{each little square can hold
at most one point}. Why? Because the two recursive calls already guarantee that any
two points on the \emph{same} side are at least $\delta$ apart, each square sits
wholly on one side, and the diagonal of a $\delta/2$-square is only
$\delta/\sqrt{2} < \delta$ -- two points in one square would be closer than
$\delta$, a contradiction. So at most $8$ points sit in the rectangle, one of
which is $p$ itself, leaving at most $7$ others to check.

\begin{center}
\begin{tikzpicture}[scale=1.5]
  \draw[step=0.5,gray!50,thin] (0,0) grid (2,1);
  \draw[thick,structurepurple] (0,0) rectangle (2,1);
  \draw[dashed,thick,trapred] (1,-0.1)--(1,1.15) node[above,font=\scriptsize]{median};
  \fill (1,0) circle (1.4pt) node[below]{$p$};
  \node[structurepurple] at (1,1.32){$2\delta \times \delta$ rectangle, $4\times2$ grid};
  \node at (0.25,0.25){\scriptsize $\tfrac{\delta}{2}$};
  \node[sourceblue,font=\scriptsize] at (0.5,0.75){left};
  \node[targetgreen,font=\scriptsize] at (1.5,0.75){right};
  \draw[<->] (2.15,0)--(2.15,1) node[midway,right]{$\delta$};
  \draw[<->] (0,-0.18)--(2,-0.18) node[midway,below]{$2\delta$};
\end{tikzpicture}
\end{center}

\begin{formalbox}
\textbf{Claim.} Let $p$ and $q$ be points in the strip with $q$ above $p$ in
$y$-order and $\mathrm{dist}(p,q) < \delta$. Then $q$ is among the next $7$ points
after $p$ in the $y$-sorted strip. \\[2pt]
\emph{Proof sketch.} Such a $q$ lies in the $2\delta\times\delta$ rectangle above
$p$, which straddles the median line: its left half is a $\delta\times\delta$
square sitting on the left side, its right half a $\delta\times\delta$ square on
the right side. Tile each half with four $\tfrac{\delta}{2}\times\tfrac{\delta}{2}$
squares, eight in all. Each little square has diameter
$\tfrac{\delta}{\sqrt2}<\delta$, and every square lies entirely on one side of the
line, so any two points inside the same square are on the same side and would be
closer than $\delta$ -- contradicting that the per-side minimum is exactly
$\delta$. Hence each square holds at most one input point, so at most $8$ points
occupy the rectangle, $p$ included, leaving at most $7$ candidates after $p$.
$\qed$
\end{formalbox}

A constant -- $7$ -- comparisons per point. That converts the strip step from
$\Th(n^2)$ into $\Th(n)$. \emph{That} is the trick.

\subsection{The algorithm and its recurrence}

\begin{algorithm}[H]
\caption{\textsc{ClosestPair}($X, Y$)}
\KwIn{Points sorted by $x$ in array $X$ and by $y$ in array $Y$ (same set).}
\KwOut{The smallest distance between any two of the points.}
\If{$\abs{X} \le 3$}{\Return brute-force minimum distance}
split $X$ at the median $x$ into left half $X_L$ and right half $X_R$\;
let $\ell$ be the median $x$-coordinate\;
partition $Y$ into $Y_L$, $Y_R$ preserving $y$-order (one linear pass)\;
$\delta_L \leftarrow \textsc{ClosestPair}(X_L, Y_L)$\;
$\delta_R \leftarrow \textsc{ClosestPair}(X_R, Y_R)$\;
$\delta \leftarrow \min(\delta_L, \delta_R)$\;
build $S \leftarrow$ points of $Y$ with $\abs{x - \ell} < \delta$ \tcp*{strip, still $y$-sorted}
\For{$i \leftarrow 1$ \KwTo $\abs{S}$}{
  \For{$j \leftarrow i+1$ \KwTo $\min(i+7, \abs{S})$}{
    $\delta \leftarrow \min(\delta, \mathrm{dist}(S[i], S[j]))$\;
  }
}
\Return $\delta$\;
\end{algorithm}

The one subtlety that makes this $\Th(n \log n)$ rather than $\Th(n \log^2 n)$:
we must \emph{not} re-sort by $y$ inside each recursive call. Instead, we pass the
already-$y$-sorted array $Y$ down and split it into $Y_L, Y_R$ in a single linear
pass each level (just like the merge step of merge sort, run in reverse). With that
care, each level of recursion does $\Th(n)$ non-recursive work.

\begin{complexitybox}
\textbf{The recurrence.} One presort costs $\Th(n \log n)$. Inside the recursion,
each call splits ($\Th(n)$), recurses on two halves, and does the strip scan with
its at-most-$7$ comparisons ($\Th(n)$). So
\[
  T(n) = 2\,T(n/2) + \Th(n).
\]
By the master theorem (Chapter~4), case~2, this solves to
$T(n) = \Th(n \log n)$. Add the one-time presort -- also $\Th(n \log n)$ -- and
the total stays $\Th(n \log n)$. We beat the $\Th(n^2)$ brute force, and we did it
with the \emph{same} recurrence as merge sort.
\end{complexitybox}

The recursion tree makes the $\Th(n \log n)$ visible: every level does $\Th(n)$
total work (the split plus the at-most-7 strip scan), and there are $\Th(\log n)$
levels.

\begin{center}
\begin{tikzpicture}[scale=0.9,
    every node/.style={font=\small},
    lvl/.style={draw,structurepurple,rounded corners,minimum width=1.5cm,inner sep=2pt}]
  \node[lvl] (a) at (0,0) {$n$};
  \node[lvl] (b) at (-2.2,-1.3) {$n/2$};
  \node[lvl] (c) at (2.2,-1.3) {$n/2$};
  \node[lvl] (d) at (-3.3,-2.6) {$n/4$};
  \node[lvl] (e) at (-1.1,-2.6) {$n/4$};
  \node[lvl] (f) at (1.1,-2.6) {$n/4$};
  \node[lvl] (g) at (3.3,-2.6) {$n/4$};
  \draw (a)--(b); \draw (a)--(c);
  \draw (b)--(d); \draw (b)--(e); \draw (c)--(f); \draw (c)--(g);
  \node at (3.3,-3.4){\vdots};   \node at (-3.3,-3.4){\vdots};
  % right-hand work column
  \node[targetgreen] at (6.0,0) {work $\Th(n)$};
  \node[targetgreen] at (6.0,-1.3) {$2\cdot\Th(n/2)=\Th(n)$};
  \node[targetgreen] at (6.0,-2.6) {$4\cdot\Th(n/4)=\Th(n)$};
  \draw[gray!50,dashed] (4.0,0.4)--(4.0,-3.0);
  \node[rotate=90,gray] at (4.3,-1.3){$\Th(\log n)$ levels};
\end{tikzpicture}
\end{center}

Summing $\Th(n)$ across $\Th(\log n)$ levels gives $\Th(n \log n)$ -- precisely the
shape of the merge-sort tree.

\begin{trapbox}
The number-one bug in closest-pair implementations is re-sorting the strip (or the
whole set) by $y$ \emph{inside} every recursive call. That sneaks an extra
$\log n$ factor in and gives $\Th(n \log^2 n)$ -- still better than $n^2$, but not
optimal, and on an exam it loses points. The fix is to thread the $y$-sorted order
through the recursion by splitting, never re-sorting. A second subtle point: use
the strip's $\delta$ \emph{from the recursive calls}, and keep shrinking it as you
find closer straddling pairs -- but always loop the at-most-7 window, not a window
that depends on the shrinking $\delta$.
\end{trapbox}

\begin{walkbox}
Convince yourself the constant really is bounded. Take the $2\delta \times \delta$
rectangle and try to cram in more than 8 points such that no two on the same side
are closer than $\delta$. You will fail -- the squares simply run out of room.
This ``packing'' style of argument, where geometry forces a constant, recurs all
over computational geometry. Where else have you seen ``the space can only hold so
many''? (Hint: hashing, and the pigeonhole principle.)
\end{walkbox}

\section{Looking back}

Step back and notice what carried the whole chapter: a single arithmetic sign.

\begin{focusbox}
\textbf{What to carry forward.}
\begin{itemize}
  \item The \emph{cross product} $u_x v_y - u_y v_x$ turns ``which way did we
        turn?'' into the sign of one number -- exact on integers, no trig, no
        division. It is the ``\texttt{<}'' of geometry.
  \item \emph{Segment intersection} is four orientation tests: each segment must
        straddle the other's line.
  \item \emph{Convex hull}: Graham's scan sorts then walks, popping right turns,
        in $\Th(n \log n)$. Jarvis's march gift-wraps in $\Th(nh)$ -- great when
        the hull is small, quadratic when it is not.
  \item \emph{Closest pair}: divide by a vertical line, recurse, and patch the
        middle strip using the seven-neighbors fact, for $\Th(n \log n)$.
  \item The deepest lesson: \emph{geometry yields to algebra}. Translate the
        picture into signs and packing arguments, and the algorithms write
        themselves.
\end{itemize}
\end{focusbox}

\begin{historybox}
For the rigorous versions of all four, read CLRS Chapter~33 next: 33.1
(line-segment properties), 33.2 (the sweep-line ``any segments intersect''),
33.3 (Graham's scan and Jarvis's march), 33.4 (closest pair). CLRS proves the
seven-neighbors bound carefully and handles every degenerate case (collinear
points, coincident points, vertical lines) that we have waved past. Having built
the intuition here, those proofs should now read as confirmations of pictures you
already hold in your head.
\end{historybox}

\section*{Exercises}
\addcontentsline{toc}{section}{Exercises}

\begin{exercise}\label{ch32:ex-orient}
Compute \textsc{Orientation}$(p_0,p_1,p_2)$ by hand for
$p_0=(1,1)$, $p_1=(4,3)$, $p_2=(2,5)$. Is the turn left, right, or straight? Then
swap $p_1$ and $p_2$ and recompute; explain why the sign flipped.
\end{exercise}

\begin{exercise}\label{ch32:ex-area}
The magnitude of the cross product is twice the area of triangle $p_0 p_1 p_2$.
Use this to find the area of the triangle with vertices $(0,0)$, $(5,0)$,
$(2,4)$. Check against the schoolbook formula $\tfrac12 \cdot \text{base} \cdot
\text{height}$.
\end{exercise}

\begin{exercise}\label{ch32:ex-seg-touch}
Do segments $\overline{p_1p_2}$ and $\overline{p_3p_4}$ intersect for
$p_1=(0,0), p_2=(4,4), p_3=(0,4), p_4=(2,2)$? Notice $p_4$ lies \emph{on} the
first segment. Walk through \textsc{SegmentsIntersect} and identify which branch
fires.
\end{exercise}

\begin{exercise}\label{ch32:ex-no-div}
Explain in two or three sentences why the orientation test is preferred over
computing and comparing slopes. Give a concrete pair of segments where the slope
approach would divide by zero but the cross product handles it cleanly.
\end{exercise}

\begin{exercise}\label{ch32:ex-graham}
Run Graham's scan by hand on the points
$(0,0),(2,0),(4,1),(2,2),(1,1),(3,1)$. Identify the anchor, the angular order, and
which points get popped. Draw the final hull.
\end{exercise}

\begin{exercise}\label{ch32:ex-jarvis-count}
For a set of $n$ points in ``convex position'' (all on the hull, e.g.\ on a
circle), exactly how many times does Jarvis's march execute its inner \texttt{for}
loop in total? Express your answer in terms of $n$ and confirm it gives
$\Th(n^2)$.
\end{exercise}

\begin{exercise}\label{ch32:ex-collinear}
In Graham's scan we pop on a non-left turn using the test
\textsc{Orientation}$(\cdot,\cdot,p_i) \le 0$. Suppose three consecutive hull
candidates are exactly collinear. With $\le 0$, is the middle one kept or popped?
What changes if you use $< 0$ instead? Describe an application that would want each
behavior.
\end{exercise}

\begin{exercise}\label{ch32:ex-strip}
In the closest-pair strip step we compare each point with the next $7$ in
$y$-order. Redo the packing argument with $\delta/2$-squares but suppose someone
claims you only need to check the next $5$ points. Construct (or argue the
impossibility of) a configuration that would break the 5-bound, and conclude why
$7$ is the safe constant.
\end{exercise}

\begin{exercise}\label{ch32:ex-recurrence}
The closest-pair recurrence is $T(n) = 2T(n/2) + \Th(n)$. Solve it from scratch
with a recursion tree: how many levels, how much work per level, and what is the
total? Then explain precisely where an accidental re-sort would change the
$\Th(n)$ term and what the new solution would be.
\end{exercise}

\begin{exercise}\label{ch32:ex-3d-challenge}
\emph{(Challenge.)} The cross product of two 3D vectors is itself a vector, not a
scalar. Look up the formula and explain how, in 3D, you would test whether four
points are coplanar (the analogue of three points being collinear). Which
arithmetic operation plays the role of our 2D sign test? \emph{(Hint: the scalar
triple product, a $3\times3$ determinant.)}
\end{exercise}
